%% kapa: exam=1 round=32 year=2021 period=2 form=A subject=appraisal_law subject_ko=감정평가관계법규 seq=37 number=37
%% source: Q-Net 감정평가사 기출문제 (cst00309) artlSeq=5208397 file=extracted/1008/5208397/A형/2021년도 제32회 감정평가사 1차 2교시 A형.pdf
%% conversion: kordoc 4.24.51 -> tools/gen1.py (auto); review=proofread-sonnet; flags=none
\begin{question}{37}
도시 및 주거환경정비법령상 정비계획 입안을 위하여 주민의견 청취절차를 거쳐야 하는 경우는? (단, 조례는 고려하지 않음)
\begin{choices}
  \item 공동이용시설 설치계획을 변경하는 경우
  \item 재난방지에 관한 계획을 변경하는 경우
  \item 정비사업시행 예정시기를 3년의 범위에서 조정하는 경우
  \item 건축물의 최고 높이를 변경하는 경우
  \item 건축물의 용적률을 20퍼센트 미만의 범위에서 확대하는 경우
\end{choices}
\end{question}
\answer{5}
